% GENERATED FILE. Edit canonical lessons, then run npm run generate:books.
% canonical-source-hash: fnv1a64:41313c63bae0e2a9
% canonical-lessons: FA-C42-shekar, FA-C42-namak, FA-C42-kare, FA-C42-angur, FA-C42-porteghal

\chapter{Sugar, Salt, Butter, Grapes, Orange}
\label{ch:fa-sugar-salt-butter-grapes-orange}
\hlchaptermodality{42}

\begin{chapteropening}
\textbf{By the end of this chapter:} \emph{I can say sugar, salt, butter, grapes and an orange.}

Complete the last lesson of chapter 42: I can say sugar, salt, butter, grapes and an orange.
\end{chapteropening}

\section[\texorpdfstring{\fa{شکر} (shekar)}{شکر (shekar)}]{\fa{شکر} (shekar) — sugar}

\label{lesson:FA-C42-shekar}



\begin{quote}
Before the new one: say the Persian for a tomato, then the Persian for a carrot.
\end{quote}

\subsection*{You'll want to know: \fa{شکر}}
\textbf{\fa{شکر}} — \emph{shekar} — \textquotedblleft{}sugar\textquotedblright{}.

\begin{cousinweb}
\emph{shekar} is usually traced to Sanskrit \emph{śarkarā}, \textquotedblleft{}grit, sugar\textquotedblright{}, through Middle Persian \emph{šakar}. Arabic \emph{sukkar} took the word from Persian, and English \emph{sugar} comes from that Arabic word by way of Latin and French.
\end{cousinweb}

\begin{grammarlens}[title={What you've built}]
One more word for this chapter.
\end{grammarlens}

\subsection*{Your turn}
\begin{itemize}
\raggedright
  \item \emph{Say it:} \emph{shekar}
  \item \emph{Say it:} \emph{shekar}, once more
  \item \emph{Say it:} say \emph{shekar}
  \item \emph{Recall:} say the Persian for a tomato, then the Persian for a carrot, then say \emph{shekar} again
\end{itemize}

\begin{tcolorbox}[breakable,colback=teal!4,colframe=teal!35!black,title={Before you move on}]
What does \fa{شکر} mean? (\textquotedblleft{}Sugar\textquotedblright{}.) Say it once more.
\end{tcolorbox}
\section[\texorpdfstring{\fa{نمک} (namak)}{نمک (namak)}]{\fa{نمک} (namak) — salt}

\label{lesson:FA-C42-namak}



\begin{quote}
Before the new one: say the Persian for a carrot, then the Persian for sugar.
\end{quote}

\subsection*{You'll want to know: \fa{نمک}}
\textbf{\fa{نمک}} — \emph{namak} — \textquotedblleft{}salt\textquotedblright{}.

\begin{cousinweb}
\emph{namak} is a native Persian word, inherited rather than borrowed: it continues Middle Persian \emph{namak}, \textquotedblleft{}salt\textquotedblright{} (EtymDB).
\end{cousinweb}

\begin{grammarlens}[title={What you've built}]
One more word for this chapter.
\end{grammarlens}

\subsection*{Your turn}
\begin{itemize}
\raggedright
  \item \emph{Say it:} \emph{namak}
  \item \emph{Say it:} \emph{namak}, once more
  \item \emph{Say it:} say \emph{namak}
  \item \emph{Recall:} say the Persian for a carrot, then the Persian for sugar, then say \emph{namak} again
\end{itemize}

\begin{tcolorbox}[breakable,colback=teal!4,colframe=teal!35!black,title={Before you move on}]
What does \fa{نمک} mean? (\textquotedblleft{}Salt\textquotedblright{}.) Say it once more.
\end{tcolorbox}
\section[\texorpdfstring{\fa{کره} (kare)}{کره (kare)}]{\fa{کره} (kare) — butter}

\label{lesson:FA-C42-kare}



\begin{quote}
Before the new one: say the Persian for sugar, then the Persian for salt.
\end{quote}

\subsection*{You'll want to know: \fa{کره}}
\textbf{\fa{کره}} — \emph{kare} — \textquotedblleft{}butter\textquotedblright{}.

\begin{cousinweb}
\emph{kare} is a native Persian word, inherited rather than borrowed: it continues Middle Persian \emph{karag}, \textquotedblleft{}butter\textquotedblright{} (EtymDB). (EtymDB also lists Arabic \emph{kura}, \textquotedblleft{}ball\textquotedblright{}, but that is a different word that looks alike.)
\end{cousinweb}

\begin{grammarlens}[title={What you've built}]
One more word for this chapter.
\end{grammarlens}

\subsection*{Your turn}
\begin{itemize}
\raggedright
  \item \emph{Say it:} \emph{kare}
  \item \emph{Say it:} \emph{kare}, once more
  \item \emph{Say it:} say \emph{kare}
  \item \emph{Recall:} say the Persian for sugar, then the Persian for salt, then say \emph{kare} again
\end{itemize}

\begin{tcolorbox}[breakable,colback=teal!4,colframe=teal!35!black,title={Before you move on}]
What does \fa{کره} mean? (\textquotedblleft{}Butter\textquotedblright{}.) Say it once more.
\end{tcolorbox}
\section[\texorpdfstring{\fa{انگور} (angur)}{انگور (angur)}]{\fa{انگور} (angur) — grapes}

\label{lesson:FA-C42-angur}



\begin{quote}
Before the new one: say the Persian for salt, then the Persian for butter.
\end{quote}

\subsection*{You'll want to know: \fa{انگور}}
\textbf{\fa{انگور}} — \emph{angur} — \textquotedblleft{}grapes\textquotedblright{}.

\begin{cousinweb}
\emph{angur} is a native Persian word, inherited rather than borrowed; it is usually traced to Middle Persian \emph{angūr}, \textquotedblleft{}grapes\textquotedblright{}.
\end{cousinweb}

\begin{grammarlens}[title={What you've built}]
One more word for this chapter.
\end{grammarlens}

\subsection*{Your turn}
\begin{itemize}
\raggedright
  \item \emph{Say it:} \emph{angur}
  \item \emph{Say it:} \emph{angur}, once more
  \item \emph{Say it:} say \emph{angur}
  \item \emph{Recall:} say the Persian for salt, then the Persian for butter, then say \emph{angur} again
\end{itemize}

\begin{tcolorbox}[breakable,colback=teal!4,colframe=teal!35!black,title={Before you move on}]
What does \fa{انگور} mean? (\textquotedblleft{}Grapes\textquotedblright{}.) Say it once more.
\end{tcolorbox}
\section[\texorpdfstring{\fa{پرتقال} (porteghâl)}{پرتقال (porteghâl)}]{\fa{پرتقال} (porteghâl) — an orange}

\label{lesson:FA-C42-porteghal}



\begin{quote}
Before the new one: say the Persian for butter, then the Persian for grapes.
\end{quote}

\subsection*{You'll want to know: \fa{پرتقال}}
\textbf{\fa{پرتقال}} — \emph{porteghâl} — \textquotedblleft{}an orange\textquotedblright{}.

\begin{cousinweb}
\emph{porteghâl} is the country name Portugal. The sweet orange is usually said to have reached the region through Portuguese traders, and Persian, like Arabic, Greek and Turkish, named the fruit after them.
\end{cousinweb}

\begin{grammarlens}[title={What you've built}]
One more word for this chapter.
\end{grammarlens}

\subsection*{Your turn}
\begin{itemize}
\raggedright
  \item \emph{Say it:} \emph{porteghâl}
  \item \emph{Say it:} \emph{porteghâl}, once more
  \item \emph{Say it:} say \emph{porteghâl}
  \item \emph{Recall:} say the Persian for butter, then the Persian for grapes, then say \emph{porteghâl} again
\end{itemize}

\begin{tcolorbox}[breakable,colback=teal!4,colframe=teal!35!black,title={Before you move on}]
What does \fa{پرتقال} mean? (\textquotedblleft{}An orange\textquotedblright{}.) Say it once more.
\end{tcolorbox}
